\section{Conclusion}
\label{sec:conclusion}

On a quarter-million-bar intraday panel, we asked what exogenous
information adds to a HAR backbone, what estimator can extract it, and
what the nonlinearity is made of. The answers form a single chain. The
exogenous panel carries real, statistically significant information ---
but the information is spread densely and weakly across hundreds of
columns, so it is invisible to selection and to unpenalized joint fits,
and recoverable only by shrinkage. The appropriate shrinkage is
block-structured: almost none on the strong low-dimensional backbone,
heavy and uniform on the wide exogenous block.
% PRODUCT BLOCK PARKED 2026-08-18: the product-block characterization of the nonlinearity and the three-block final model are held out of the rendered paper; the two sentences below are the two-block rewrite, plus the options reading.
\iffalse
The recoverable
nonlinearity is a sparse, frozen set of pairwise products gated on the
session clock, distillable from tree fits into one hundred explicit
columns, and valuable only when shrunk harder still and refit often. The
final model is ridge with a block-diagonal Tikhonov penalty, linear in its parameters throughout:
three nested blocks, three fixed penalties, one exact rolling solve, and it outperforms the
benchmark by a margin an order of magnitude beyond the fourth decimal
place at strong paired significance.
\fi
% 2026-10-06 parked (main results re-scored on the plain back-transform, user decision):
% The recoverable nonlinearity
% is small beside that information effect and is concentrated at pairwise
% order --- products of columns the design already contains, gated on the
% session clock --- a specification statement the tree fits deliver. The
% final model is ridge with a block-diagonal Tikhonov penalty, linear in
% its parameters throughout: two nested blocks, two fixed penalties, one
% exact rolling solve, and it outperforms the benchmark by a margin an
% order of magnitude beyond the fourth decimal place at strong paired
% significance.
The recoverable nonlinearity is not small: on the wide design the hedged
tree bank improves on the two-block ridge by \unifTreeBankHedgeVsBlkTwoDQ{} at
DM \unifTreeBankHedgeVsBlkTwoDM{}, more than the ridge improves on the
benchmark, and a decomposition of a tree fit places most of its explainable
variance in additive terms carried by the HAR block. The linear model is
ridge with a block-diagonal Tikhonov penalty, linear in its parameters
throughout: two nested blocks and one exact rolling solve. With its two
penalties fixed it improves on the benchmark by \unifBlkTwoUserDelta{} at DM
\unifBlkTwoUserDM{}, short of conventional significance; with each block's
penalty re-selected causally, by \unifBlkTwoTunedDelta{} at DM
\unifBlkTwoTunedDM{}.
% PARKED 2026-08-28: rest-of-session / multi-horizon held out of the rendered paper.
\iffalse
Extended to the variance remaining in the session, the
same per-bar forecast serves as a multi-horizon forecast at every bar
using only information available then, and the two-block ridge's
advantage over the incumbent carries to that horizon.
\fi
% D1 2026-09-29 (change 11): re-read on Section 5.4 as it now stands; the session-bar
% version is kept, commented, below.  No number is quoted here.
% 2026-10-06 parked (recalibration removed, user decision):
% The paper closes by reading the forecasts against the last thirty minutes of
% the expiration-day chain (Section~\ref{sec:close_option_results}). Under one
% recalibration and on the same days, the sign of the gap between forecast and
% quoted implied variance carries a direction edge on the straddle: every
% forecast improves on always short in point estimate, and a ridge fitted for
% the 15:30--16:00 bar alone, on inputs a forecaster can rebuild live, is
% resolved above both always short and the paper's own forecast at both fills.
% The forecast built for the whole panel is therefore not the one to trade at
% its last bar --- but the accuracy gain behind the difference is not
% significant, no single step from one forecast to the other is resolved alone
% (the largest is fitting the traded bar on its own), and the edge is a tail
% trade: a few long days carry it, the sample holds few of them, and the
% trade's live form has shown no edge after the sample ends. Sizing on the size of the gap does not improve
% on its sign.
The paper closes by reading the forecasts against the last thirty minutes of
the expiration-day chain (Section~\ref{sec:close_option_results}). With every
forecast used as its plain back-transform and on the same days, the sign of the
gap between forecast and quoted implied variance carries a direction edge on
% 2026-10-06 parked (no headline model; live-feasible removed from the paper; user decisions):
% the straddle: each of the paper's eight forecasts improves on always short, and a
% ridge fitted for the 15:30--16:00 bar alone, on inputs a forecaster can rebuild live,
% trades above both always short and the paper's own forecast at both fills ---
% significantly only against always short.
% The forecast built for the whole panel need not be the one to trade at
% its last bar --- but the ridge fitted for that bar is no more accurate than the
% paper's forecast, no single step from one forecast to the other moves the trade
% significantly (the largest is fitting the traded bar on its own), the Sharpe
% ratios of the many per-bar forecasts are not an established order, and the edge is a tail
% trade: a few long days carry it, the sample holds few of them, and the
% trade's live form has shown no edge after the sample ends. Sizing on the size of the gap does not improve
% on its sign.
the straddle: each of the paper's eight forecasts improves on always short. Across the
many forecasts of the closing study's master table, including those fitted for the
15:30--16:00 bar alone, the differences in the trade's Sharpe ratio are mostly not
significant and do not order the models, no single step from one forecast to another
moves the trade significantly, and a better forecast of the bar's variance is not by
% 2026-10-06 parked (running appendix re-run on the plain forecast: D.5 held out, sizing re-run):
% itself a better trade. The trade's live form has shown no edge after the sample ends.
% Sizing on the size of the gap does not improve on its sign.
itself a better trade. Whether sizing on the size of the gap improves on its sign turns
on the sample, in point estimate; the trade's live form after the sample ends is not yet
scored with the plain forecast.
% PARKED 2026-09-29 (D1 change 11): the session-bar version.  Verbatim:
% The paper closes by reading that one-bar forecast against the last
% thirty minutes of the expiration-day chain
% (Section~\ref{sec:close_option_results}): the sign of the gap between
% forecast and quoted implied is the tradable content --- a $\pm 1$
% portfolio on $\mathrm{sign}(s)$ improves on the always-short control for
% every one of the eight forecast columns, the block-diagonal ridge of record
% among them, and survives a crossed spread, though only one column's paired
% improvement clears zero on a bootstrap interval and the columns cannot be
% ordered against one another --- while sizing on the magnitude of that gap
% does not improve on its sign. The ridge of record is fitted on the design
% carrying the FOMC calendar block; four comparators are still on the earlier
% design, and the table reports a diagnostic row that isolates what those
% calendar channels are worth on this trade.
% OPTIONS PARKED 2026-08-18: option-market reading held out of the rendered paper.
\iffalse
Read against the same-day index-option market through a
delta-hedged near-the-money portfolio, its remaining-variance forecast is
compared with the incumbent's under one fixed rule; the reading is a
demonstration under stated constants.
\fi

% 2026-10-06 parked (main results re-scored on the plain back-transform, user decision):
% Two findings sit beside the model rather than inside it, and may matter
% more broadly. First, the evaluation protocol is a modeling choice: the
% causal calibrated back-transform beats the retrospective Duan convention
% on every arm, and its adaptive second moment --- not any conditional-mean
% device --- is among the largest measured effects in the study. Second,
% causal hyperparameter tuning on short validation tails is repeatedly
% worse than well-chosen constants, because the price of a selection
% mistake is asymmetric across estimator families. Both findings are about
% the discipline surrounding the model, and both are measured, not
% assumed.
One finding sits beside the model rather than inside it, and may matter
more broadly: what causal hyperparameter tuning on short validation tails
buys depends on the estimator. Weighting the tree experts beats selecting
one (\unifIncrTreeHedgeDQ{} at DM \unifIncrTreeHedgeDM{}), and re-selecting
the two block penalties causally improves on fixing them. It is measured,
not assumed.

% 2026-10-06 parked (main results re-scored on the plain back-transform, user decision):
% The limitations are stated where they bind. Results are computed on one
% asset complex and one loss; levels are comparable only within one
% preprocessing regime and one protocol, and arms that survive only under
% the legacy protocol are labeled as such in the appendix. The
% conditional-variance panel remains on the legacy scalar back-transform
% and has not been re-scored under the protocol used in the body.
% The natural extensions are equally concrete: promoting the
% mean-reverting second-moment refinement to the protocol (a re-score of
% every comparison);
The limitations are stated where they bind. Results are computed on one
asset complex and one loss, and levels are comparable only within one
preprocessing regime. The natural extension is equally concrete:
% PRODUCT BLOCK PARKED 2026-08-18: the tree-expert-bank extension is stated only as a comparison against the product block, so it is held out; the list is now two items.
\iffalse
the frozen tree-expert bank, which puts the trees and
the product block on identical footing;
\fi
% 2026-10-06 parked (main results re-scored on the plain back-transform, user decision):
% and model confidence sets over
% the feature sets once per-bar losses are persisted. None of these
% changes the chain above; each tightens a link in it.
model confidence sets over
the feature sets once per-bar losses are persisted. It would not change
the chain above; it would tighten a link in it.
